\documentclass[11pt]{article}
\usepackage[a4paper,margin=0.65in]{geometry}
\usepackage[T1]{fontenc}
\usepackage{amsmath,amssymb,amsfonts}
\usepackage{graphicx}
\usepackage{pdfpages}
\usepackage[english,shorthands=off]{babel}
\usepackage{fancyhdr}
\usepackage[bookmarks=false,colorlinks=false]{hyperref}
\usepackage[protrusion=true,expansion=false]{microtype}
\emergencystretch=3em
\tolerance=600
\providecommand{\modd}{\mathrm{mod}}
\providecommand{\Disc}{\mathrm{Disc}}
\providecommand{\canont}{}
\providecommand{\asterism}{*}
\providecommand{\Hessian}{H}
\providecommand{\tome}{}
\providecommand{\page}{p.}
\pagestyle{fancy}\fancyhf{}\fancyhead[L]{Pages 388--400}\fancyhead[R]{Cayley --- Collected Papers, Vol.~IX}\renewcommand{\headrulewidth}{0.4pt}
\setlength{\headheight}{15pt}

\begin{document}

%% ========================================================================
%% PAPER [607] -- A Memoir on Prepotentials (continued, scan pp. 368--380)
%%                PDF pp. 388--400
%% ========================================================================
\section*{\textbf{607.} A Memoir on Prepotentials (continued)}

\textit{[Continuation of the same paper; articles 75 (concluded)--95 of the Annexes.]}

\bigskip

\noindent\textbf{[p.~368]}\quad We hence have
\[
\resizebox{\textwidth}{!}{$\displaystyle
\text{Ring-integral}
= \frac{f^{s-1}}{(f^{2}-\kappa^{2})^{q+\frac{1}{2}}}\,\frac{\Gamma\tfrac{1}{2}s\,\Gamma\tfrac{1}{2}s}{\Gamma(\tfrac{1}{2}s+q)\,\Gamma(\tfrac{1}{2}-q)}\int_{0}^{\infty} t^{\frac{1}{2}s+q-1}(t+f^{2}-\kappa^{2})^{-q-\frac{1}{2}}(t+f^{2})^{-\frac{1}{2}s-q+\frac{1}{2}}\,dt$}
\]
\[
\resizebox{\textwidth}{!}{$\displaystyle
= \frac{f^{s-1}}{(f^{2}-\kappa^{2})^{q+\frac{1}{2}}}\,\frac{\Gamma\tfrac{1}{2}s\,\Gamma\tfrac{1}{2}s}{\Gamma(\tfrac{1}{2}s+q)\,\Gamma(\tfrac{1}{2}-q)}\int_{0}^{\infty} t^{q-\frac{1}{2}}(t+f^{2}-\kappa^{2})^{\frac{1}{2}s-q-1}(t+f^{2})^{-\frac{1}{2}s-q}\,dt$}
\]
\[
\resizebox{\textwidth}{!}{$\displaystyle
= f^{s-1}\,\frac{\Gamma\tfrac{1}{2}s\,\Gamma\tfrac{1}{2}s}{\Gamma(\tfrac{1}{2}-q)\,\Gamma(\tfrac{1}{2}s+q)}\int_{0}^{\infty} t^{-q-\frac{1}{2}}(t+f^{2}-\kappa^{2})^{\frac{1}{2}s+q-1}(t+f^{2})^{-\frac{1}{2}s+q}\,dt$}
\]
\[
\resizebox{\textwidth}{!}{$\displaystyle
= f^{s-1}\,\frac{\Gamma\tfrac{1}{2}s\,\Gamma\tfrac{1}{2}s}{\Gamma(\tfrac{1}{2}s-q)\,\Gamma(\tfrac{1}{2}+q)}\int_{0}^{\infty} t^{\frac{1}{2}s-q-1}(t+f^{2}-\kappa^{2})^{q-\frac{1}{2}}(t+f^{2})^{-\frac{1}{2}-q}\,dt.$}
\]

As a verification write $\kappa=0$, the four integrals are
\[
\int_{0}^{\infty}\frac{t^{\frac{1}{2}s+q-1}\,dt}{(t+f^{2})^{q+\frac{1}{2}}} = f^{s-q}\,\frac{\Gamma(\tfrac{1}{2}s+q)\,\Gamma(\tfrac{1}{2}-q)}{\Gamma(\tfrac{1}{2}s+\tfrac{1}{2})},
\]
\[
\int_{0}^{\infty}\frac{t^{q-\frac{1}{2}}\,dt}{(t+f^{2})^{q+\frac{1}{2}}} = f^{-s-q}\,\frac{\Gamma(\tfrac{1}{2}+q)\,\Gamma(\tfrac{1}{2}s-q)}{\Gamma(\tfrac{1}{2}s+\tfrac{1}{2})},
\]
\[
\int_{0}^{\infty}\frac{t^{-q-\frac{1}{2}}\,dt}{(t+f^{2})^{\frac{1}{2}-q}} = f^{-s-q}\,\frac{\Gamma(\tfrac{1}{2}-q)\,\Gamma(\tfrac{1}{2}s+q)}{\Gamma(\tfrac{1}{2}s+\tfrac{1}{2})},
\]
\[
\int_{0}^{\infty}\frac{t^{\frac{1}{2}s-q-1}\,dt}{(t+f^{2})^{\frac{1}{2}+q}} = f^{-s-q}\,\frac{\Gamma(\tfrac{1}{2}s-q)\,\Gamma(\tfrac{1}{2}+q)}{\Gamma(\tfrac{1}{2}s+\tfrac{1}{2})};
\]
hence from each of them
\[
\text{Ring-integral} = \frac{1}{f^{2q}}\,\frac{\Gamma\tfrac{1}{2}s\,\Gamma\tfrac{1}{2}s}{\Gamma(\tfrac{1}{2}s+\tfrac{1}{2})},
\]
which is, in fact, the value obtained from
\[
\text{Ring-integral} = \frac{2^{s-1}f^{s}}{(f+\kappa)^{s+2q}}\,\Pi\!\left(\tfrac{1}{2}s,\,\tfrac{1}{2}s,\,\tfrac{1}{2}s+q,\,\tfrac{4\kappa f}{(f+\kappa)^{2}}\right)
\]
on putting therein $\kappa=0$; viz.\ the value is
\[
= \frac{2^{s-1}f^{s}}{f^{s+2q}}\,\Pi(\tfrac{1}{2}s,\,\tfrac{1}{2}s,\,\tfrac{1}{2}s+q,\,0).
\]

\bigskip

\noindent\textbf{76.}\quad For an external point $\kappa>f$, $\sqrt{1-u} = \dfrac{\kappa-f}{\kappa+f}$, and therefore $v=\dfrac{f^{2}}{\kappa^{2}}$.
\begin{align*}
\Pi(\tfrac{1}{2}s,\tfrac{1}{2}s,\tfrac{1}{2}s+q,u)
&= \left(\frac{\kappa+f}{\kappa}\right)^{s+2q} 2^{-q}\,\frac{\Gamma\tfrac{1}{2}s\,\Gamma\tfrac{1}{2}s}{\Gamma(\tfrac{1}{2}s+q)\,\Gamma(\tfrac{1}{2}-q)}\,\Pi\!\left(\tfrac{1}{2}s+q,\,\tfrac{1}{2}-q,\,\tfrac{1}{2}+q,\,\tfrac{f^{2}}{\kappa^{2}}\right) \\[3pt]
&= \left(\frac{\kappa+f}{\kappa}\right)^{s+2q} 2^{-q}\,\frac{\Gamma\tfrac{1}{2}s\,\Gamma\tfrac{1}{2}s}{\Gamma(\tfrac{1}{2}+q)\,\Gamma(\tfrac{1}{2}s-q)}\,\Pi\!\left(\tfrac{1}{2}+q,\,\tfrac{1}{2}s-q,\,\tfrac{1}{2}s+q,\,\tfrac{f^{2}}{\kappa^{2}}\right) \\[3pt]
&= \left(\frac{\kappa+f}{\kappa}\right)^{s}\!\left(\frac{\kappa-f}{\kappa}\right)^{-2q} 2^{-q}\,\frac{\Gamma\tfrac{1}{2}s\,\Gamma\tfrac{1}{2}s}{\Gamma(\tfrac{1}{2}-q)\,\Gamma(\tfrac{1}{2}s+q)}\,\Pi\!\left(\tfrac{1}{2}-q,\,\tfrac{1}{2}s+q,\,\tfrac{1}{2}s-q,\,\tfrac{f^{2}}{\kappa^{2}}\right) \\[3pt]
&= \left(\frac{\kappa+f}{\kappa}\right)^{s}\!\left(\frac{\kappa-f}{\kappa}\right)^{-2q} 2^{-q}\,\frac{\Gamma\tfrac{1}{2}s\,\Gamma\tfrac{1}{2}s}{\Gamma(\tfrac{1}{2}s-q)\,\Gamma(\tfrac{1}{2}+q)}\,\Pi\!\left(\tfrac{1}{2}s-q,\,\tfrac{1}{2}+q,\,\tfrac{1}{2}-q,\,\tfrac{f^{2}}{\kappa^{2}}\right);
\end{align*}

\bigskip

\noindent\textbf{[p.~369]}\quad where the $\Pi$-functions on the right hand are respectively
\[
\int_{0}^{1}\frac{x^{\frac{1}{2}s+q-1}(1-x)^{-q-\frac{1}{2}}\,dx}{(\kappa^{2}-f^{2}x)^{q+\frac{1}{2}}}
= \kappa^{-2q+1}\int_{0}^{\infty}\frac{t^{\frac{1}{2}s+q-1}\,dt}{(t+\kappa^{2}-f^{2})^{q+\frac{1}{2}}(t+\kappa^{2})^{\frac{1}{2}s+q-\frac{1}{2}}},
\]
\[
\int_{0}^{1}\frac{x^{q-\frac{1}{2}}(1-x)^{\frac{1}{2}s-q-1}\,dx}{(\kappa^{2}-f^{2}x)^{\frac{1}{2}s+q}}
= \kappa^{-2q+1}\int_{0}^{\infty}\frac{t^{q-\frac{1}{2}}(t+\kappa^{2}-f^{2})^{\frac{1}{2}s-q-1}\,dt}{(t+\kappa^{2})^{\frac{1}{2}s+q}},
\]
\[
\int_{0}^{1}\frac{x^{-q-\frac{1}{2}}(1-x)^{\frac{1}{2}s+q-1}\,dx}{(\kappa^{2}-f^{2}x)^{\frac{1}{2}s-q}}
= \kappa^{-2q+1}\int_{0}^{\infty}\frac{t^{-q-\frac{1}{2}}(t+\kappa^{2}-f^{2})^{\frac{1}{2}s+q-1}\,dt}{(t+\kappa^{2})^{\frac{1}{2}s-q}},
\]
\[
\int_{0}^{1}\frac{x^{\frac{1}{2}s-q-1}(1-x)^{q-\frac{1}{2}}\,dx}{(\kappa^{2}-f^{2}x)^{\frac{1}{2}-q}}
= \kappa^{-2q+1}\int_{0}^{\infty}\frac{t^{\frac{1}{2}s-q-1}(t+\kappa^{2}-f^{2})^{q-\frac{1}{2}}\,dt}{(t+\kappa^{2})^{\frac{1}{2}+q}}.
\]

We have then
\[
\resizebox{\textwidth}{!}{$\displaystyle
\text{Ring-integral}
= \frac{f^{s-1}}{(\kappa^{2}-f^{2})^{q+\frac{1}{2}}}\,\frac{\Gamma\tfrac{1}{2}s\,\Gamma\tfrac{1}{2}s}{\Gamma(\tfrac{1}{2}s+q)\,\Gamma(\tfrac{1}{2}-q)}\int_{0}^{\infty} t^{\frac{1}{2}s+q-1}(t+\kappa^{2}-f^{2})^{-q-\frac{1}{2}}(t+\kappa^{2})^{-\frac{1}{2}s-q+\frac{1}{2}}\,dt$}
\]
\[
\resizebox{\textwidth}{!}{$\displaystyle
= \frac{f^{s-1}}{(\kappa^{2}-f^{2})^{q+\frac{1}{2}}}\,\frac{\Gamma\tfrac{1}{2}s\,\Gamma\tfrac{1}{2}s}{\Gamma(\tfrac{1}{2}+q)\,\Gamma(\tfrac{1}{2}s-q)}\int_{0}^{\infty} t^{q-\frac{1}{2}}(t+\kappa^{2}-f^{2})^{\frac{1}{2}s-q-1}(t+\kappa^{2})^{-\frac{1}{2}s-q}\,dt$}
\]
\[
\resizebox{\textwidth}{!}{$\displaystyle
= f^{s-1}\,\frac{\Gamma\tfrac{1}{2}s\,\Gamma\tfrac{1}{2}s}{\Gamma(\tfrac{1}{2}-q)\,\Gamma(\tfrac{1}{2}s+q)}\int_{0}^{\infty} t^{-q-\frac{1}{2}}(t+\kappa^{2}-f^{2})^{\frac{1}{2}s+q-1}(t+\kappa^{2})^{-\frac{1}{2}s+q}\,dt$}
\]
\[
\resizebox{\textwidth}{!}{$\displaystyle
= f^{s-1}\,\frac{\Gamma\tfrac{1}{2}s\,\Gamma\tfrac{1}{2}s}{\Gamma(\tfrac{1}{2}s-q)\,\Gamma(\tfrac{1}{2}+q)}\int_{0}^{\infty} t^{\frac{1}{2}s-q-1}(t+\kappa^{2}-f^{2})^{q-\frac{1}{2}}(t+\kappa^{2})^{-\frac{1}{2}-q}\,dt.$}
\]

Observe that in II.\ and III.\ the integrals, except as to the limits, are the same as in the corresponding formul{\ae} for the interior point.

If in the integrals we put $t+\kappa^{2}-f^{2}$ in place of $t$, and ultimately suppose $\kappa$ indefinitely large in comparison with $f$, they severally become
\[
\int_{-\kappa^{2}+f^{2}}^{\infty}\frac{(t+\kappa^{2}-f^{2})^{\frac{1}{2}s+q-1}\,dt}{(t+\kappa^{2})^{q+\frac{1}{2}}} = \kappa^{s-1}\,\frac{\Gamma(\tfrac{1}{2}s+q)\,\Gamma(\tfrac{1}{2}-q)}{\Gamma(\tfrac{1}{2}s+\tfrac{1}{2})},
\]
\[
\int_{-\kappa^{2}+f^{2}}^{\infty}\frac{(t+\kappa^{2}-f^{2})^{q-\frac{1}{2}}\,t^{\frac{1}{2}s-q-1}\,dt}{(t+\kappa^{2})^{\frac{1}{2}s+q}} = \kappa^{s-1}\,\frac{\Gamma(\tfrac{1}{2}+q)\,\Gamma(\tfrac{1}{2}s-q)}{\Gamma(\tfrac{1}{2}s+\tfrac{1}{2})},
\]
\[
\int_{-\kappa^{2}+f^{2}}^{\infty}\frac{(t+\kappa^{2}-f^{2})^{-q-\frac{1}{2}}\,t^{\frac{1}{2}s+q-1}\,dt}{(t+\kappa^{2})^{\frac{1}{2}s-q}} = \kappa^{s-1}\,\frac{\Gamma(\tfrac{1}{2}-q)\,\Gamma(\tfrac{1}{2}s+q)}{\Gamma(\tfrac{1}{2}s+\tfrac{1}{2})},
\]
\[
\int_{-\kappa^{2}+f^{2}}^{\infty}\frac{(t+\kappa^{2}-f^{2})^{\frac{1}{2}s-q-1}\,t^{q-\frac{1}{2}}\,dt}{(t+\kappa^{2})^{\frac{1}{2}+q}} = \kappa^{s-1}\,\frac{\Gamma(\tfrac{1}{2}s-q)\,\Gamma(\tfrac{1}{2}+q)}{\Gamma(\tfrac{1}{2}s+\tfrac{1}{2})};
\]
and they all four give
\[
\text{Ring-integral} = \frac{f^{s-1}}{\kappa^{2q+1}}\cdot\frac{\Gamma\tfrac{1}{2}s\,\Gamma\tfrac{1}{2}s}{\Gamma(\tfrac{1}{2}s+\tfrac{1}{2})},
\]
which agrees with the value
\[
= \frac{2^{s-1}f^{s}}{(\kappa+f)^{s+2q}}\,\Pi(\tfrac{1}{2}s,\,\tfrac{1}{2}s,\,\tfrac{1}{2}s+q,\,0),
\]
when $\tfrac{f}{\kappa}$ is indefinitely large.

\bigskip
\bigskip

\noindent\textbf{[p.~370]}\quad \textbf{77.}\quad We come now to the disk-integral,
\[
\int\!\!\int\frac{y^{s-1}\,dx\,dy}{\{(x-\kappa)^{2}+y^{2}\}^{\frac{1}{2}s+q}},
\]
over the circle $x^{2}+y^{2}=f^{2}$. Writing $x=\kappa+\rho\cos\phi$, $y=\rho\sin\phi$, we have $dx\,dy=\rho\,d\rho\,d\phi$, and the integral therefore is
\[
\int\!\!\int\frac{\sin^{s-1}\phi\,d\rho\,d\phi}{\rho^{2q}},
\]
where the integration in regard to $\rho$ is performed at once; viz.\ the integral is
\[
= \frac{1}{1-2q}\int(\rho^{1-2q})\sin^{s-1}\phi\,d\phi\,;
\]
or multiplying by $2$, in order that the integration may be taken only over the semicircle, $y=$ positive, this is
\[
= \frac{1}{\tfrac{1}{2}-q}\int(\rho^{1-2q})\sin^{s-1}\phi\,d\phi,
\]
the term $(\rho^{1-2q})$ being taken between the proper limits.

\bigskip

\noindent\textbf{78.}\quad Consider first an interior point $\kappa<f$. As already mentioned, we exclude an indefinitely small circle radius $\epsilon$, and the limits for $\rho$ are from $\rho=\epsilon$ to $\rho=$ its value at

\bigskip

\textit{[Figure: a circle, with $O$ on the horizontal axis as origin, an interior point $P$ near $O$, and the line $OP$ extended through a point on the circle; $y$-axis vertical.]}

\bigskip

\noindent the circumference; viz.\ if here $x=f\cos\theta$, $y=f\sin\theta$, then we have $f\cos\theta=\kappa+\rho\cos\phi$, $f\sin\theta=\rho\sin\phi$, and consequently
\[
\rho^{2} = \kappa^{2}+f^{2}-2\kappa f\cos\theta,
\]
\[
\sin\phi = \frac{f\sin\theta}{\rho},\qquad \cos\phi = \frac{f\cos\theta-\kappa}{\rho} = \frac{f\sin\theta}{\sqrt{\kappa^{2}+f^{2}-2\kappa f\cos\theta}},
\]
and the integral therefore is
\[
= \frac{1}{\tfrac{1}{2}-q}\int\!\left\{\frac{f^{s-1}\sin^{s-1}\theta}{(\kappa^{2}+f^{2}-2\kappa f\cos\theta)^{\frac{1}{2}s+q-\frac{1}{2}}} - \epsilon^{1-2q}\sin^{s-1}\phi\right\}d\phi.
\]

As regards the second term, this is $\displaystyle =\frac{\epsilon^{1-2q}}{\tfrac{1}{2}-q}\int\sin^{s-1}\phi\,d\phi$, from $\phi=0$ to $\phi=\pi$; or, what is the same thing, we may multiply by $2$ and take the integral from $\phi=0$ to $\phi=\tfrac{\pi}{2}$.

\bigskip

\noindent\textbf{[p.~371]}\quad Writing then $\sin\phi=\sqrt{x}$, and consequently $\sin^{s-1}\phi\,d\phi=\tfrac{1}{2}x^{\frac{1}{2}s-1}(1-x)^{-\frac{1}{2}}\,dx$, the term is
\[
= -\frac{\epsilon^{1-2q}}{\tfrac{1}{2}-q}\,\frac{\Gamma\tfrac{1}{2}s\,\Gamma\tfrac{1}{2}}{\Gamma(\tfrac{1}{2}s+\tfrac{1}{2})};
\]
the value of the disk-integral is
\[
= \frac{f^{s-1}}{\tfrac{1}{2}-q}\int\!\frac{\sin^{s-1}\theta\,d\theta}{(\kappa^{2}+f^{2}-2\kappa f\cos\theta)^{\frac{1}{2}s+q-\frac{1}{2}}} - \frac{\epsilon^{1-2q}}{\tfrac{1}{2}-q}\,\frac{\Gamma\tfrac{1}{2}s\,\Gamma\tfrac{1}{2}}{\Gamma(\tfrac{1}{2}s+\tfrac{1}{2})}.
\]
But we have
\[
\sin\phi = \frac{f\sin\theta}{\rho},\qquad \cos\phi = \frac{f\cos\theta-\kappa}{\rho},
\]
and thence
\[
\tan\phi = \frac{f\sin\theta}{f\cos\theta-\kappa},\qquad \sec^{2}\phi\,d\phi = \frac{f(f-\kappa\cos\theta)\,d\theta}{(f\cos\theta-\kappa)^{2}};
\]
that is,
\[
d\phi = \frac{f(f-\kappa\cos\theta)\,d\theta}{\rho^{2}} = \frac{f(f-\kappa\cos\theta)\,d\theta}{f^{2}+\kappa^{2}-2\kappa f\cos\theta},
\]
or, what is the same thing,
\[
= \tfrac{1}{2}\,\frac{\{(f^{2}-\kappa^{2}) + (f^{2}+\kappa^{2}-2\kappa f\cos\theta)\}\,d\theta}{f^{2}+\kappa^{2}-2\kappa f\cos\theta};
\]
the expression for the disk-integral is therefore
\[
\resizebox{\textwidth}{!}{$\displaystyle
= \frac{\tfrac{1}{2}f^{s-1}}{\tfrac{1}{2}-q}\int_{0}^{\pi}\!\frac{\sin^{s-1}\theta\,\{(f^{2}-\kappa^{2})+(f^{2}+\kappa^{2}-2\kappa f\cos\theta)\}\,d\theta}{(f^{2}+\kappa^{2}-2\kappa f\cos\theta)^{\frac{1}{2}s+q}} - \frac{\epsilon^{1-2q}}{\tfrac{1}{2}-q}\,\frac{\Gamma\tfrac{1}{2}s\,\Gamma\tfrac{1}{2}}{\Gamma(\tfrac{1}{2}s+\tfrac{1}{2})}.$}
\]

\bigskip

\noindent\textbf{79.}\quad Writing as before $\cos\tfrac{1}{2}\theta=\sqrt{x}$, $\sin\tfrac{1}{2}\theta=\sqrt{1+x}$, \&c., and $u=\dfrac{4\kappa f}{(\kappa+f)^{2}}$, this is
\[
\resizebox{\textwidth}{!}{$\displaystyle
= \frac{2^{s-1}f^{s-1}}{(\tfrac{1}{2}-q)(\kappa+f)^{s+2q-1}}\!\left\{\frac{f^{2}-\kappa^{2}}{(\kappa+f)^{2}}\,\Pi(\tfrac{1}{2}s,\,\tfrac{1}{2}s,\,\tfrac{1}{2}s+q,\,u) + \Pi(\tfrac{1}{2}s,\,\tfrac{1}{2}s,\,\tfrac{1}{2}s+q-1,\,u)\right\} - \frac{\epsilon^{1-2q}}{\tfrac{1}{2}-q}\,\frac{\Gamma\tfrac{1}{2}s\,\Gamma\tfrac{1}{2}}{\Gamma(\tfrac{1}{2}s+\tfrac{1}{2})}.$}
\]

As a verification, observe that, if $\kappa=0$, each of the $\Pi$-functions becomes
\[
= \int_{0}^{1} x^{\frac{1}{2}s-1}(1-x)^{-\frac{1}{2}}\,dx = \frac{\Gamma\tfrac{1}{2}s\,\Gamma\tfrac{1}{2}}{\Gamma(\tfrac{1}{2}s+\tfrac{1}{2})};
\]
hence the whole first term is $=\dfrac{f^{s-1}\cdot f}{\tfrac{1}{2}-q}\cdot\dfrac{\Gamma\tfrac{1}{2}s\,\Gamma\tfrac{1}{2}}{\Gamma(\tfrac{1}{2}s+\tfrac{1}{2})}$, viz.\ this is $\dfrac{f^{s-2q}}{\tfrac{1}{2}-q}\cdot\dfrac{\Gamma\tfrac{1}{2}s\,\Gamma\tfrac{1}{2}}{\Gamma(\tfrac{1}{2}s+\tfrac{1}{2})}$, and the complete value is
\[
= \frac{1}{\tfrac{1}{2}-q}\,\frac{\Gamma\tfrac{1}{2}s\,\Gamma\tfrac{1}{2}}{\Gamma(\tfrac{1}{2}s+\tfrac{1}{2})}\,(f^{s-2q}-\epsilon^{s-2q}),
\]
vanishing, as it should do, if $f=\epsilon$.

\bigskip

\noindent\textbf{80.}\quad In the case of an exterior point $\kappa>f$, the process is somewhat different; but the result is of a like form. We have
\[
\text{Disk-integral} = \frac{1}{\tfrac{1}{2}-q}\int(\rho_{1}^{1-2q}-\rho^{1-2q})\sin^{s-1}\phi\,d\phi,
\]

\bigskip

\noindent\textbf{[p.~372]}\quad where $\rho_{1}$ refers to the point $M'$ and $\rho$ to the point $M$. Attending first to the integral $\displaystyle\int \rho^{1-2q}\sin^{s-1}\phi\,d\phi$, and writing as before $f\cos\theta=\kappa+\rho\cos\phi$, $f\sin\theta=\rho\sin\phi$, this is
\[
= f^{s-1}\int\!\frac{\sin^{s-1}\theta\,d\theta}{(\kappa^{2}+f^{2}-2\kappa f\cos\theta)^{\frac{1}{2}s+q-\frac{1}{2}}}
\]
\[
= \tfrac{1}{2}f^{s-1}\int\!\frac{\sin^{s-1}\theta\,\{-(\kappa^{2}-f^{2})+(f^{2}+\kappa^{2}-2\kappa f\cos\theta)\}\,d\theta}{(f^{2}+\kappa^{2}-2\kappa f\cos\theta)^{\frac{1}{2}s+q}},
\]
the inferior and the superior limits being here the values of $\theta$ which correspond to the points $N$, $A$ respectively, say $\theta+\alpha$, and $\theta=0$; hence, reversing the sign and interchanging the two limits, the value of $-\!\displaystyle\int\rho^{1-2q}\sin^{s-1}\phi\,d\phi$ is the above integral taken from $0$ to $\alpha$. But similarly the value of $+\!\displaystyle\int\rho_{1}^{1-2q}\sin^{s-1}\phi\,d\phi$ is the same integral taken from $\alpha$ to $\pi$. For the two terms together, the value is the same integral from $0$ to $\pi$; viz.\ we thus find
\[
\text{Disk-integral} = \tfrac{1}{2}\,\frac{f^{s-1}}{\tfrac{1}{2}-q}\int_{0}^{\pi}\!\frac{\sin^{s-1}\theta\,\{-(\kappa^{2}-f^{2})+(f^{2}+\kappa^{2}-2\kappa f\cos\theta)\}\,d\theta}{(f^{2}+\kappa^{2}-2\kappa f\cos\theta)^{\frac{1}{2}s+q}};
\]
or, writing as before $\cos\tfrac{1}{2}\theta=\sqrt{x}$, \&c., and $u=\dfrac{4\kappa f}{(\kappa+f)^{2}}$, this is
\[
\resizebox{\textwidth}{!}{$\displaystyle
= \frac{2^{s-1}f^{s-1}}{(\tfrac{1}{2}-q)(\kappa+f)^{s+2q-1}}\!\left\{-\frac{\kappa^{2}-f^{2}}{(\kappa+f)^{2}}\,\Pi(\tfrac{1}{2}s,\,\tfrac{1}{2}s,\,\tfrac{1}{2}s+q,\,u) + \Pi(\tfrac{1}{2}s,\,\tfrac{1}{2}s,\,\tfrac{1}{2}s+q-1,\,u)\right\}.$}
\]

\bigskip

\noindent\textbf{81.}\quad As a verification, suppose that $\kappa$ is indefinitely large: we must recur to the last preceding formula; the value is thus
\[
= \frac{f^{s}}{(\tfrac{1}{2}-q)\,\kappa^{s+2q-1}}\int_{0}^{\pi}\!\frac{\sin^{s-1}\theta\left(-\cos\theta+\dfrac{f}{\kappa}\right)}{\left(1-\dfrac{2f}{\kappa}\cos\theta\right)^{\frac{1}{2}s+q}}\,d\theta,
\]
viz.\ this is
\[
= \frac{f^{s}}{(\tfrac{1}{2}-q)\,\kappa^{s+2q-1}}\int_{0}^{\pi}\sin^{s-1}\theta\left\{-\cos\theta+[1-(s+2q)\cos^{2}\theta]\,\frac{f}{\kappa}\right\}d\theta,
\]
where the integral of the first term vanishes; the value is thus
\[
= \frac{f^{s+1}}{(\tfrac{1}{2}-q)\,\kappa^{s+2q}}\int_{0}^{\pi}\sin^{s-1}\theta\,[1-(s+2q)\cos^{2}\theta]\,d\theta,
\]

\bigskip

\noindent\textbf{[p.~373]}\quad where we may multiply by $2$ and take the integral from $0$ to $\tfrac{\pi}{2}$. Writing $\sin\theta=\sqrt{x}$, the value is
\[
= \frac{f^{s+1}}{(\tfrac{1}{2}-q)\,\kappa^{s+2q}}\int_{0}^{1} x^{\frac{1}{2}s-1}\,[1-(s+2q)(1-x)]\,(1-x)^{-\frac{1}{2}}\,dx,
\]
where the integral is
\[
= \frac{\Gamma\tfrac{1}{2}s\,\Gamma\tfrac{1}{2}}{\Gamma(\tfrac{1}{2}s+\tfrac{1}{2})}\!\left\{1 - \frac{\tfrac{1}{2}(s+2q)}{\tfrac{1}{2}s+\tfrac{1}{2}}\right\}
= \frac{\Gamma\tfrac{1}{2}s\,\Gamma\tfrac{1}{2}}{\Gamma(\tfrac{1}{2}s+\tfrac{1}{2})}\,\frac{\tfrac{1}{2}-q}{\tfrac{1}{2}s+\tfrac{1}{2}},
\]
and hence the value is
\[
= \frac{f^{s+1}}{\kappa^{s+2q}}\,\frac{\Gamma\tfrac{1}{2}s\,\Gamma\tfrac{1}{2}}{\Gamma(\tfrac{1}{2}s+\tfrac{3}{2})};
\]
viz.\ this is $\displaystyle =\frac{1}{\kappa^{s+2q}}\int y^{s-1}\,dx\,dy$, over the circle $x^{2}+y^{2}=f^{2}$, as is easily verified.

\bigskip

\noindent\textbf{82.}\quad Reverting to the interior point $\kappa<f$.
\[
\resizebox{\textwidth}{!}{$\displaystyle
\text{Disk-integral} = \frac{2^{s-1}f^{s-1}}{(\tfrac{1}{2}-q)(\kappa+f)^{s+2q-1}}\!\left\{\frac{f^{2}-\kappa^{2}}{(\kappa+f)^{2}}\,\Pi(\tfrac{1}{2}s,\,\tfrac{1}{2}s,\,\tfrac{1}{2}s+q,\,u) + \Pi(\tfrac{1}{2}s,\,\tfrac{1}{2}s,\,\tfrac{1}{2}s+q-1,\,u)\right\} - \frac{\epsilon^{1-2q}}{\tfrac{1}{2}-q}\,\frac{\Gamma\tfrac{1}{2}s\,\Gamma\tfrac{1}{2}}{\Gamma(\tfrac{1}{2}s+\tfrac{1}{2})};$}
\]
then reducing the expression in $\{\ \}$ by the transformations for $\Pi(\tfrac{1}{2}s,\tfrac{1}{2}s,\tfrac{1}{2}s+q,u)$ and the like transformations for $\Pi(\tfrac{1}{2}s,\tfrac{1}{2}s,\tfrac{1}{2}s+q-1,u)$, the term in $\{\ \}$ may be expressed in the four forms:
\[
\resizebox{\textwidth}{!}{$\displaystyle
1.\quad \frac{2^{-q}\,\Gamma\tfrac{1}{2}s\,\Gamma\tfrac{1}{2}s}{\Gamma(\tfrac{1}{2}s+q)\,\Gamma(\tfrac{1}{2}-q)}\,\frac{(f+\kappa)^{s+2q-1}}{f^{s+2q-1}}\!\left[\left(1-\tfrac{\kappa^{2}}{f^{2}}\right)\!\Pi(\tfrac{1}{2}s+q,\,\tfrac{1}{2}-q,\,\tfrac{1}{2}+q,\tfrac{\kappa^{2}}{f^{2}})+\tfrac{\tfrac{1}{2}+q-\tfrac{1}{2}}{\tfrac{1}{2}-q}\Pi(\tfrac{1}{2}s+q-1,\,\tfrac{3}{2}-q,\,-\tfrac{1}{2}+q,\tfrac{\kappa^{2}}{f^{2}})\right],$}
\]
\[
\resizebox{\textwidth}{!}{$\displaystyle
2.\quad \frac{2^{-q}\,\Gamma\tfrac{1}{2}s\,\Gamma\tfrac{1}{2}s}{\Gamma(\tfrac{1}{2}+q)\,\Gamma(\tfrac{1}{2}s-q)}\,\frac{(f+\kappa)^{s+2q-1}}{f^{s+2q-1}}\!\left[\left(1-\tfrac{\kappa^{2}}{f^{2}}\right)\!\Pi(\tfrac{1}{2}+q,\,\tfrac{1}{2}s-q,\,\tfrac{1}{2}s+q,\tfrac{\kappa^{2}}{f^{2}})+\tfrac{\tfrac{1}{2}+q-\tfrac{1}{2}}{\tfrac{1}{2}-q}\Pi(-\tfrac{1}{2}+q,\,\tfrac{1}{2}s-q+1,\,\tfrac{1}{2}s+q-1,\tfrac{\kappa^{2}}{f^{2}})\right],$}
\]
\[
\resizebox{\textwidth}{!}{$\displaystyle
3.\quad \frac{2^{-q}\,\Gamma\tfrac{1}{2}s\,\Gamma\tfrac{1}{2}s}{\Gamma(\tfrac{1}{2}-q)\,\Gamma(\tfrac{1}{2}s+q)}\,\frac{(f+\kappa)^{s-1}(f-\kappa)^{-2q}}{f^{s-1}}\!\left[\left(1-\tfrac{\kappa^{2}}{f^{2}}\right)\!\Pi(\tfrac{1}{2}-q,\,\tfrac{1}{2}s+q,\,\tfrac{1}{2}s-q,\tfrac{\kappa^{2}}{f^{2}})+\tfrac{\tfrac{1}{2}+q-\tfrac{1}{2}}{\tfrac{1}{2}-q}\Pi(\tfrac{1}{2}-q,\,\tfrac{1}{2}s+q,\,\tfrac{1}{2}s-q,\tfrac{\kappa^{2}}{f^{2}})\right],$}
\]
\[
\resizebox{\textwidth}{!}{$\displaystyle
4.\quad \frac{2^{-q}\,\Gamma\tfrac{1}{2}s\,\Gamma\tfrac{1}{2}s}{\Gamma(\tfrac{1}{2}s-q)\,\Gamma(\tfrac{1}{2}+q)}\,\frac{(f+\kappa)^{s-1}(f-\kappa)^{-2q}}{f^{s-1}}\!\left[\left(1-\tfrac{\kappa^{2}}{f^{2}}\right)\!\Pi(\tfrac{1}{2}s-q,\,\tfrac{1}{2}+q,\,\tfrac{1}{2}-q,\tfrac{\kappa^{2}}{f^{2}})+\tfrac{\tfrac{1}{2}+q-\tfrac{1}{2}}{\tfrac{1}{2}-q}\Pi(\tfrac{1}{2}s-q+1,\,-\tfrac{1}{2}+q,\,\tfrac{3}{2}-q,\tfrac{\kappa^{2}}{f^{2}})\right].$}
\]

\bigskip

\noindent\textbf{[p.~374]}\quad \textbf{83.}\quad The first and the fourth of these are susceptible of a reduction which does not appear to be applicable to the second and the third. Consider in general the function
\[
(1-v)\,\Pi(\alpha,\,\beta,\,1-\beta,\,v) + \frac{\alpha-1}{\beta}\,\Pi(\alpha-1,\,\beta+1,\,-\beta,\,v);
\]
the second $\Pi$-function is here
\[
\int_{0}^{1} x^{\alpha-2}(1-x\,.\,1-vx)^{\beta}\,dx;
\]
viz.\ this is
\[
= \frac{x^{\alpha-1}}{\alpha-1}(1-x\,.\,1-vx)^{\beta} - \frac{1}{\alpha-1}\int_{0}^{1} x^{\alpha-1}\frac{d}{dx}(1-x\,.\,1-vx)^{\beta}\,dx,
\]
or, since the first term vanishes between the limits, this is
\[
= \frac{\beta}{\alpha-1}\int x^{\alpha-1}(1-x\,.\,1-vx)^{\beta-1}(1+v-2vx)\,dx,
\]
\[
= \frac{\beta}{\alpha-1}\!\left\{(1+v)\,\Pi(\alpha,\,\beta,\,1-\beta,\,v) - 2v\!\int_{0}^{1} x^{\alpha}(1-x\,.\,1-vx)^{\beta-1}\,dx\right\}.
\]
Hence the two $\Pi$-functions together are
\[
= (1-v+1+v)\!\int_{0}^{1} x^{\alpha-1}(1-x\,.\,1-vx)^{\beta-1}\,dx - 2\!\int_{0}^{1} vx\,.\,x^{\alpha-1}(1-x\,.\,1-vx)^{\beta-1}\,dx,
\]
\[
= 2\!\int_{0}^{1} x^{\alpha-1}(1-x)^{\beta-1}(1-vx)^{\beta}\,dx,
\]
that is,
\[
(1-v)\,\Pi(\alpha,\,\beta,\,1-\beta,\,v) + \tfrac{\alpha-1}{\beta}\,\Pi(\alpha-1,\,\beta+1,\,-\beta,\,v) = 2\,\Pi(\alpha,\,\beta,\,-\beta,\,v).
\]
We have therefore
\[
\resizebox{\textwidth}{!}{$\displaystyle
\left(1-\tfrac{\kappa^{2}}{f^{2}}\right)\!\Pi(\tfrac{1}{2}s+q,\,\tfrac{1}{2}-q,\,\tfrac{1}{2}+q,\,\tfrac{\kappa^{2}}{f^{2}}) + \tfrac{\tfrac{1}{2}s+q-1}{\tfrac{1}{2}-q}\,\Pi(\tfrac{1}{2}s+q-1,\,\tfrac{3}{2}-q,\,-\tfrac{1}{2}+q,\,\tfrac{\kappa^{2}}{f^{2}}) = 2\,\Pi(\tfrac{1}{2}s+q,\,\tfrac{1}{2}-q,\,-\tfrac{1}{2}+q,\,\tfrac{\kappa^{2}}{f^{2}});$}
\]
and from the same equation, written in the form
\[
\Pi(\alpha-1,\,\beta+1,\,-\beta,\,v) + \tfrac{\beta}{\alpha-1}(1-v)\,\Pi(\alpha,\,\beta,\,1-\beta,\,v) = 2\tfrac{\beta}{\alpha-1}\,\Pi(\alpha,\,\beta,\,-\beta,\,v),
\]
we obtain
\[
\resizebox{\textwidth}{!}{$\displaystyle
\Pi(\tfrac{1}{2}s-q+1,\,-\tfrac{1}{2}+q,\,\tfrac{3}{2}-q,\,v) + (1-v)\,\tfrac{-\tfrac{1}{2}+q}{\tfrac{1}{2}s-q}\,\Pi(\tfrac{1}{2}s-q,\,\tfrac{1}{2}+q,\,\tfrac{1}{2}-q,\,v) = 2\tfrac{-\tfrac{1}{2}+q}{\tfrac{1}{2}s-q}\,\Pi(\tfrac{1}{2}s-q,\,\tfrac{1}{2}+q,\,-\tfrac{1}{2}-q,\,v).$}
\]

\bigskip

\noindent\textbf{84.}\quad Hence the terms in $[\ ]$ in the first and the fourth expressions in No.\ 82 are
\[
= \frac{2^{q+1}\,\Gamma\tfrac{1}{2}s\,\Gamma\tfrac{1}{2}s}{\Gamma(\tfrac{1}{2}s+q)\,\Gamma(\tfrac{1}{2}-q)}\,\frac{(f+\kappa)^{s+2q-1}}{f^{s+2q-1}}\,\Pi(\tfrac{1}{2}s+q,\,\tfrac{1}{2}-q,\,-\tfrac{1}{2}+q,\,\tfrac{\kappa^{2}}{f^{2}}),
\]
\[
= \frac{2^{-q}(-\tfrac{1}{2}+q)\,\Gamma\tfrac{1}{2}s\,\Gamma\tfrac{1}{2}s}{\Gamma(\tfrac{1}{2}s-q+1)\,\Gamma(\tfrac{1}{2}+q)}\,\frac{(f+\kappa)^{s-1}(f-\kappa)^{-2q}}{f^{s-1}}\,\Pi(\tfrac{1}{2}s-q+1,\,-\tfrac{1}{2}+q,\,\tfrac{1}{2}-q,\,\tfrac{\kappa^{2}}{f^{2}}),
\]

\bigskip

\noindent\textbf{[p.~375]}\quad respectively; the corresponding values of the disk-integral are
\[
\resizebox{\textwidth}{!}{$\displaystyle
\frac{\Gamma\tfrac{1}{2}s\,\Gamma\tfrac{1}{2}s}{\Gamma(\tfrac{1}{2}-q)\,\Gamma(\tfrac{1}{2}s+q)}\,f^{s-2q}\,\Pi(\tfrac{1}{2}s+q,\,\tfrac{1}{2}-q,\,-\tfrac{1}{2}+q,\,\tfrac{\kappa^{2}}{f^{2}}) - \frac{\epsilon^{1-2q}}{\tfrac{1}{2}-q}\,\frac{\Gamma\tfrac{1}{2}s\,\Gamma\tfrac{1}{2}}{\Gamma(\tfrac{1}{2}s+\tfrac{1}{2})},$}
\]
\[
\resizebox{\textwidth}{!}{$\displaystyle
\frac{-\Gamma\tfrac{1}{2}s\,\Gamma\tfrac{1}{2}s}{\Gamma(\tfrac{1}{2}s-q+1)\,\Gamma(\tfrac{1}{2}+q)}\!\left(\frac{f-\kappa}{f}\right)^{1-2q}\!\Pi(\tfrac{1}{2}s-q+1,\,-\tfrac{1}{2}+q,\,\tfrac{1}{2}-q,\,\tfrac{\kappa^{2}}{f^{2}}) - \frac{\epsilon^{1-2q}}{\tfrac{1}{2}-q}\,\frac{\Gamma\tfrac{1}{2}s\,\Gamma\tfrac{1}{2}}{\Gamma(\tfrac{1}{2}s+\tfrac{1}{2})},$}
\]
which we may again verify by writing therein $\kappa=0$, viz.\ the $\Pi$-functions thus become
\[
\frac{\Gamma(\tfrac{1}{2}s+q)\,\Gamma(\tfrac{1}{2}-q)}{\Gamma(\tfrac{1}{2}s+\tfrac{1}{2})} \quad\text{and}\quad \frac{\Gamma(\tfrac{1}{2}s-q+1)\,\Gamma(-\tfrac{1}{2}+q)}{\Gamma(\tfrac{1}{2}s+\tfrac{1}{2})},
\]
and consequently the integral is
\[
= \frac{1}{\tfrac{1}{2}-q}\,\frac{\Gamma\tfrac{1}{2}s\,\Gamma\tfrac{1}{2}}{\Gamma(\tfrac{1}{2}s+\tfrac{1}{2})}\,(f^{s-2q}-\epsilon^{s-2q}).
\]

\bigskip

\noindent\textbf{85.}\quad But the forms nevertheless belong to a system of four. In the formulae
\begin{align*}
\Pi(\alpha,\,\beta,\,\gamma,\,v)
&= \frac{\Gamma\alpha\,\Gamma\beta}{\Gamma\gamma\,\Gamma(\alpha+\beta-\gamma)}\,\Pi(\gamma,\,\alpha+\beta-\gamma,\,\alpha,\,v) \\[3pt]
&= (1-v)^{\gamma-\beta-\gamma}\,\Pi(\beta,\,\alpha,\,\alpha+\beta-\gamma,\,v) \\[3pt]
&= (1-v)^{\gamma-\beta-\gamma}\,\frac{\Gamma\alpha\,\Gamma\beta}{\Gamma(\alpha+\beta-\gamma)\,\Gamma\gamma}\,\Pi(\alpha+\beta-\gamma,\,\gamma,\,\beta,\,v),
\end{align*}
writing $\alpha=\tfrac{1}{2}s+q$, $\beta=\tfrac{1}{2}-q$, $\gamma=-\tfrac{1}{2}+q$, we deduce
\begin{align*}
\Pi(\tfrac{1}{2}s+q,\,\tfrac{1}{2}-q,\,-\tfrac{1}{2}+q,\,v)
&= \frac{\Gamma(\tfrac{1}{2}s+q)\,\Gamma(\tfrac{1}{2}-q)}{\Gamma(-\tfrac{1}{2}+q)\,\Gamma(\tfrac{1}{2}s-q+1)}\,\Pi(-\tfrac{1}{2}+q,\,\tfrac{1}{2}s-q+1,\,\tfrac{1}{2}s+q,\,v) \\[3pt]
&= (1-v)^{1-s}\,\Pi(\tfrac{1}{2}-q,\,\tfrac{1}{2}s+q,\,\tfrac{1}{2}s-q+1,\,v) \\[3pt]
&= (1-v)^{1-s}\,\frac{\Gamma(\tfrac{1}{2}s+q)\,\Gamma(\tfrac{1}{2}-q)}{\Gamma(\tfrac{1}{2}s-q+1)\,\Gamma(-\tfrac{1}{2}+q)}\,\Pi(\tfrac{1}{2}s-q+1,\,-\tfrac{1}{2}+q,\,\tfrac{1}{2}-q,\,v);
\end{align*}
and the last-mentioned values of the disk-integral may thus be written in the four forms
\[
\frac{\Gamma\tfrac{1}{2}s\,\Gamma\tfrac{1}{2}s}{\Gamma(\tfrac{1}{2}-q)\,\Gamma(\tfrac{1}{2}s+q)}\,f^{s-2q}\,\Pi(\tfrac{1}{2}s+q,\,\tfrac{1}{2}-q,\,-\tfrac{1}{2}+q,\,\tfrac{\kappa^{2}}{f^{2}}) \quad-\text{term in } \epsilon,
\]
\[
\frac{-\Gamma\tfrac{1}{2}s\,\Gamma\tfrac{1}{2}s}{\Gamma(\tfrac{1}{2}+q)\,\Gamma(\tfrac{1}{2}s-q+1)}\,f^{s-2q}\,\Pi(-\tfrac{1}{2}+q,\,\tfrac{1}{2}s-q+1,\,\tfrac{1}{2}s+q,\,\tfrac{\kappa^{2}}{f^{2}}) \quad-\quad ,, \quad,
\]
\[
\frac{\Gamma\tfrac{1}{2}s\,\Gamma\tfrac{1}{2}s}{\Gamma(\tfrac{1}{2}-q)\,\Gamma(\tfrac{1}{2}s+q)}\!\left(\frac{f-\kappa}{f}\right)^{1-s}\!\Pi(\tfrac{1}{2}-q,\,\tfrac{1}{2}s+q,\,\tfrac{1}{2}s-q+1,\,\tfrac{\kappa^{2}}{f^{2}}) \quad-\quad ,, \quad,
\]
\[
\frac{-\Gamma\tfrac{1}{2}s\,\Gamma\tfrac{1}{2}s}{\Gamma(\tfrac{1}{2}+q)\,\Gamma(\tfrac{1}{2}s-q+1)}\!\left(\frac{f-\kappa}{f}\right)^{1-s}\!\Pi(\tfrac{1}{2}s-q+1,\,-\tfrac{1}{2}+q,\,\tfrac{1}{2}-q,\,\tfrac{\kappa^{2}}{f^{2}}) \quad-\quad ,, \quad;
\]

\bigskip

\noindent\textbf{[p.~376]}\quad and since the last of these is in fact the second of the original forms, it is clear that, if instead of the first we had taken the second of the original forms, we should have obtained again the same system of four forms.

\bigskip

\noindent\textbf{86.}\quad Writing as before $x=\dfrac{t}{t+f^{2}-\kappa^{2}}$, \&c., the forms are
\[
\resizebox{\textwidth}{!}{$\displaystyle
\frac{\Gamma\tfrac{1}{2}s\,\Gamma\tfrac{1}{2}s}{\Gamma(\tfrac{1}{2}-q)\,\Gamma(\tfrac{1}{2}s+q)}\,(f^{2}-\kappa^{2})^{q-\frac{1}{2}}\!\int_{0}^{\infty} t^{\frac{1}{2}s+q-1}(t+f^{2}-\kappa^{2})^{-q-\frac{1}{2}}(t+f^{2})^{-\frac{1}{2}s-q+\frac{1}{2}}\,dt \quad-\text{term in } \epsilon,$}
\]
\[
\resizebox{\textwidth}{!}{$\displaystyle
\frac{-\Gamma\tfrac{1}{2}s\,\Gamma\tfrac{1}{2}s}{\Gamma(\tfrac{1}{2}+q)\,\Gamma(\tfrac{1}{2}s-q+1)}\,f^{s+1}(f^{2}-\kappa^{2})^{-q-\frac{1}{2}}\!\int_{0}^{\infty} t^{-q-\frac{1}{2}}(t+f^{2}-\kappa^{2})^{\frac{1}{2}s-q}(t+f^{2})^{-\frac{1}{2}s-q-1}\,dt \quad-\quad ,, \quad,$}
\]
\[
\resizebox{\textwidth}{!}{$\displaystyle
\frac{\Gamma\tfrac{1}{2}s\,\Gamma\tfrac{1}{2}s}{\Gamma(\tfrac{1}{2}-q)\,\Gamma(\tfrac{1}{2}s+q)}\,f^{s-1}(f^{2}-\kappa^{2})^{-q-\frac{1}{2}}\!\int_{0}^{\infty} t^{\frac{1}{2}s-q-1}(t+f^{2}-\kappa^{2})^{q-\frac{1}{2}}(t+f^{2})^{-\frac{1}{2}s+q}\,dt \quad-\text{term in } \epsilon,$}
\]
\[
\resizebox{\textwidth}{!}{$\displaystyle
\frac{-\Gamma\tfrac{1}{2}s\,\Gamma\tfrac{1}{2}s}{\Gamma(\tfrac{1}{2}+q)\,\Gamma(\tfrac{1}{2}s-q+1)}\,(f^{2}-\kappa^{2})^{q+\frac{1}{2}}\!\int_{0}^{\infty} t^{\frac{1}{2}s-q-1}(t+f^{2}-\kappa^{2})^{-q-\frac{1}{2}}(t+f^{2})^{-\frac{1}{2}-q}\,dt \quad-\quad ,, \quad.$}
\]

\bigskip

\noindent\textbf{87.}\quad The third of these possesses a remarkable property. Write $mf$ instead of $f$, and at the same time change $t$ into $m^{2}t$: the integral becomes
\[
\resizebox{\textwidth}{!}{$\displaystyle
\frac{\Gamma\tfrac{1}{2}s\,\Gamma\tfrac{1}{2}s}{\Gamma(\tfrac{1}{2}-q)\,\Gamma(\tfrac{1}{2}s+q)}\,m^{s-1}f^{s-1}\!\int_{0}^{\infty} t^{\frac{1}{2}s-q-1}\{m^{2}(t+f^{2})-\kappa^{2}\}^{q-\frac{1}{2}}\,m^{s+1-2q}(t+f^{2})^{-\frac{1}{2}s+q-\frac{1}{2}}\,dt \quad-\text{term in } \epsilon;$}
\]
and hence, writing $mf=f+\delta f$ or $m=1+\dfrac{\delta f}{f}$, and therefore $m^{2}=1+\dfrac{2\delta f}{f}$, the value is
\[
\resizebox{\textwidth}{!}{$\displaystyle
\frac{\Gamma\tfrac{1}{2}s\,\Gamma\tfrac{1}{2}s}{\Gamma(\tfrac{1}{2}-q)\,\Gamma(\tfrac{1}{2}s+q)}\,f^{s-1}\!\int_{0}^{\infty} t^{\frac{1}{2}s-q-1}\!\left\{(t+f^{2}-\kappa^{2}) + \frac{2\delta f}{f}(t+f^{2})\right\}^{q-\frac{1}{2}}\!(t+f^{2})^{-\frac{1}{2}s-q+\frac{1}{2}}\,dt \quad-\text{term in } \epsilon.$}
\]
Hence the term in $\delta f$ is
\[
\resizebox{\textwidth}{!}{$\displaystyle
= 2(-q+\tfrac{1}{2})\,\frac{\delta f}{f}\,\frac{\Gamma\tfrac{1}{2}s\,\Gamma\tfrac{1}{2}s}{\Gamma(\tfrac{1}{2}-q)\,\Gamma(\tfrac{1}{2}s+q)}\,f^{s-1}\!\int_{0}^{\infty} t^{\frac{1}{2}s-q-1}(t+f^{2}-\kappa^{2})^{q-\frac{3}{2}}(t+f^{2})^{-\frac{1}{2}s+q+\frac{1}{2}}\,dt,$}
\]
\[
\resizebox{\textwidth}{!}{$\displaystyle
= \delta f \text{ into expression}\quad \frac{\Gamma\tfrac{1}{2}s\,\Gamma\tfrac{1}{2}s}{\Gamma(\tfrac{1}{2}-q)\,\Gamma(\tfrac{1}{2}s+q)}\,f^{s-2}\!\int_{0}^{\infty} t^{\frac{1}{2}s-q-1}(t+f^{2}-\kappa^{2})^{q-\frac{3}{2}}(t+f^{2})^{-\frac{1}{2}s+q+\frac{1}{2}}\,dt,$}
\]
where the factor which multiplies $\delta f$ is, as it should be, the ring-integral; it in fact agrees with one of the expressions previously obtained for this integral.

\bigskip

\noindent\textbf{88.}\quad Similarly for an exterior point $\kappa>f$; starting in like manner from, Disk-integral
\[
\resizebox{\textwidth}{!}{$\displaystyle
= \frac{2^{s-1}f^{s-1}}{(\tfrac{1}{2}-q)(\kappa+f)^{s+2q-1}}\!\left\{-\frac{\kappa^{2}-f^{2}}{(\kappa+f)^{2}}\,\Pi(\tfrac{1}{2}s,\,\tfrac{1}{2}s,\,\tfrac{1}{2}s+q,\,u) + \Pi(\tfrac{1}{2}s,\,\tfrac{1}{2}s,\,\tfrac{1}{2}s+q-1,\,u)\right\},$}
\]
and reducing in like manner, the term in $\{\ \}$ may be expressed in the four forms
\[
\resizebox{\textwidth}{!}{$\displaystyle
1.\quad \frac{2^{-q}\,\Gamma\tfrac{1}{2}s\,\Gamma\tfrac{1}{2}s}{\Gamma(\tfrac{1}{2}+q)\,\Gamma(\tfrac{1}{2}s-q)}\,\frac{(\kappa+f)^{s+2q-1}}{\kappa^{s+2q-1}}\!\left[-\!\left(1-\tfrac{f^{2}}{\kappa^{2}}\right)\!\Pi(\tfrac{1}{2}+q,\,\tfrac{1}{2}s-q,\,\tfrac{1}{2}s+q,\,\tfrac{f^{2}}{\kappa^{2}}) + \tfrac{-\tfrac{1}{2}+q}{\tfrac{1}{2}+q}\,\Pi(-\tfrac{1}{2}+q,\,\tfrac{1}{2}s-q+1,\,\tfrac{1}{2}s+q-1,\,\tfrac{f^{2}}{\kappa^{2}})\right],$}
\]

\bigskip

\noindent\textbf{[p.~377]}\quad
\[
\resizebox{\textwidth}{!}{$\displaystyle
2.\quad \frac{\Gamma\tfrac{1}{2}s\,\Gamma\tfrac{1}{2}s}{\Gamma(\tfrac{1}{2}s+q)\,\Gamma(\tfrac{1}{2}-q)}\,\frac{(\kappa+f)^{s+2q-1}}{\kappa^{s+2q-1}}\!\left[-\!\left(1-\tfrac{f^{2}}{\kappa^{2}}\right)\!\Pi(\tfrac{1}{2}s+q,\,\tfrac{1}{2}-q,\,\tfrac{1}{2}+q,\,\tfrac{f^{2}}{\kappa^{2}}) + \tfrac{\tfrac{1}{2}+q-1}{\tfrac{1}{2}-q}\,\Pi(\tfrac{1}{2}s+q-1,\,\tfrac{3}{2}-q,\,-\tfrac{1}{2}+q,\,\tfrac{f^{2}}{\kappa^{2}})\right],$}
\]
\[
\resizebox{\textwidth}{!}{$\displaystyle
3.\quad \frac{\Gamma\tfrac{1}{2}s\,\Gamma\tfrac{1}{2}s}{\Gamma(\tfrac{1}{2}-q)\,\Gamma(\tfrac{1}{2}s+q)}\,\frac{(\kappa+f)^{s-1}(\kappa-f)^{-2q}}{\kappa^{s-1}}\!\left[-\Pi(\tfrac{1}{2}-q,\,\tfrac{1}{2}s+q,\,\tfrac{1}{2}s-q,\,\tfrac{f^{2}}{\kappa^{2}}) + \!\left(1-\tfrac{f^{2}}{\kappa^{2}}\right)\!\tfrac{\tfrac{1}{2}+q-1}{\tfrac{1}{2}-q}\,\Pi(\tfrac{1}{2}-q,\,\tfrac{1}{2}s+q,\,\tfrac{1}{2}s-q,\,\tfrac{f^{2}}{\kappa^{2}})\right],$}
\]
\[
\resizebox{\textwidth}{!}{$\displaystyle
4.\quad \frac{\Gamma\tfrac{1}{2}s\,\Gamma\tfrac{1}{2}s}{\Gamma(\tfrac{1}{2}s-q+1)\,\Gamma(\tfrac{1}{2}+q)}\,\frac{(\kappa+f)^{s-1}(\kappa-f)^{-2q}}{\kappa^{s-1}}\!\left[-\Pi(\tfrac{1}{2}s-q+1,\,-\tfrac{1}{2}+q,\,\tfrac{3}{2}-q,\,\tfrac{f^{2}}{\kappa^{2}}) + \!\left(1-\tfrac{f^{2}}{\kappa^{2}}\right)\!\tfrac{-\tfrac{1}{2}+q}{\tfrac{1}{2}+q}\,\Pi(\tfrac{1}{2}s-q+1,\,-\tfrac{1}{2}+q,\,\tfrac{3}{2}-q,\,\tfrac{f^{2}}{\kappa^{2}})\right].$}
\]

\bigskip

\noindent\textbf{89.}\quad For the reduction of the first and the fourth of these, we have to consider
\[
-(1-v)\,\Pi(\alpha,\,\beta,\,1-\beta,\,v) + \tfrac{\alpha-1}{\beta}\,\Pi(\alpha-1,\,\beta+1,\,-\beta,\,v);
\]
viz.\ this is
\[
(-1+v+1+v)\!\int_{0}^{1} x^{\alpha-1}(1-x\,.\,1-vx)^{\beta-1}\,dx - 2\!\int_{0}^{1} vx\,.\,x^{\alpha-1}(1-x\,.\,1-vx)^{\beta-1}\,dx,
\]
\[
= 2v\!\int_{0}^{1} x^{\alpha-1}(1-x)(1-x\,.\,1-vx)^{\beta-1}\,dx,
\]
\[
= 2v\,\Pi(\alpha,\,\beta+1,\,-\beta+1,\,v);
\]
that is,
\[
-(1-v)\,\Pi(\alpha,\,\beta,\,1-\beta,\,v) + \tfrac{\alpha-1}{\beta}\,\Pi(\alpha-1,\,\beta+1,\,-\beta,\,v) = 2v\,\Pi(\alpha,\,\beta+1,\,-\beta+1,\,v).
\]
I repeat, for comparison, the foregoing equation
\[
+(1-v)\,\Pi(\alpha,\,\beta,\,1-\beta,\,v) + \tfrac{\alpha-1}{\beta}\,\Pi(\alpha-1,\,\beta+1,\,-\beta,\,v) = 2\,\Pi(\alpha,\,\beta,\,-\beta,\,v);
\]
by adding and subtracting these we obtain two new formulae; for reduction of the fourth formula, the equation may be written
\[
-\Pi(\alpha-1,\,\beta+1,\,-\beta,\,v) + (1-v)\,\tfrac{\beta}{\alpha-1}\,\Pi(\alpha,\,\beta,\,1-\beta,\,v) = -2\tfrac{\beta}{\alpha-1}v\,\Pi(\alpha,\,\beta+1,\,-\beta+1,\,v).
\]

\bigskip

\noindent\textbf{90.}\quad But it is sufficient to consider the first formula; the term in $[\ ]$ is
\[
= \frac{2^{q+1}\,\Gamma\tfrac{1}{2}s\,\Gamma\tfrac{1}{2}s}{\Gamma(\tfrac{1}{2}s+q)\,\Gamma(\tfrac{1}{2}-q)}\!\left(\frac{\kappa+f}{\kappa}\right)^{s+2q-1}\!\left(\frac{\kappa-f}{\kappa}\right)^{-2q}\!\Pi(\tfrac{1}{2}+q,\,\tfrac{1}{2}s-q+1,\,\tfrac{1}{2}s+q,\,\tfrac{f^{2}}{\kappa^{2}});
\]
and the corresponding value of the disk-integral is
\[
= \frac{\Gamma\tfrac{1}{2}s\,\Gamma\tfrac{1}{2}s}{\Gamma(\tfrac{1}{2}+q)\,\Gamma(\tfrac{1}{2}-q)}\,\frac{f^{s+1}}{\kappa^{s+2q}}\,\Pi(\tfrac{1}{2}+q,\,\tfrac{1}{2}s-q+1,\,\tfrac{1}{2}s+q,\,\tfrac{f^{2}}{\kappa^{2}}).
\]

\bigskip

\noindent\textbf{[p.~378]}\quad which we may again verify by taking therein $\kappa$ indefinitely large; viz.\ the value is then $\displaystyle =\frac{f^{s+1}}{\kappa^{s+2q}}\,\frac{\Gamma\tfrac{1}{2}s\,\Gamma\tfrac{1}{2}}{\Gamma(\tfrac{1}{2}s+\tfrac{3}{2})}$, as above. It is the first of a system of four forms, the others of which are
\[
\resizebox{\textwidth}{!}{$\displaystyle
\frac{\Gamma\tfrac{1}{2}s\,\Gamma\tfrac{1}{2}s}{\Gamma(\tfrac{1}{2}s+q)\,\Gamma(\tfrac{1}{2}-q)}\,\frac{f^{s+1}}{\kappa^{s+2q}}\!\left(\frac{\kappa-f}{\kappa}\right)^{1-s}\!\Pi(\tfrac{1}{2}+q,\,\tfrac{1}{2}s-q+1,\,\tfrac{1}{2}s+q,\,\tfrac{f^{2}}{\kappa^{2}}),$}
\]
\[
\resizebox{\textwidth}{!}{$\displaystyle
\frac{\Gamma\tfrac{1}{2}s\,\Gamma\tfrac{1}{2}s}{\Gamma(\tfrac{1}{2}-q)\,\Gamma(\tfrac{1}{2}s+q)}\,\frac{f^{s+1}}{\kappa^{s+2q}}\!\left(\frac{\kappa-f}{\kappa}\right)^{1-s}\!\Pi(\tfrac{1}{2}-q,\,\tfrac{1}{2}s+q,\,\tfrac{1}{2}s-q+1,\,\tfrac{f^{2}}{\kappa^{2}}),$}
\]
\[
\resizebox{\textwidth}{!}{$\displaystyle
\frac{\Gamma\tfrac{1}{2}s\,\Gamma\tfrac{1}{2}s}{\Gamma(\tfrac{1}{2}s-q+1)\,\Gamma(\tfrac{1}{2}+q)}\,\frac{f^{s+1}}{\kappa^{s+2q}}\!\left(\frac{\kappa-f}{\kappa}\right)^{1-s}\!\Pi(\tfrac{1}{2}s-q+1,\,-\tfrac{1}{2}+q,\,\tfrac{3}{2}-q,\,\tfrac{f^{2}}{\kappa^{2}}).$}
\]
And hence, writing as before $x=\dfrac{t}{t+f^{2}-\kappa^{2}}$, \&c., the four values are
\[
\resizebox{\textwidth}{!}{$\displaystyle
\frac{\Gamma\tfrac{1}{2}s\,\Gamma\tfrac{1}{2}s}{\Gamma(\tfrac{1}{2}+q)\,\Gamma(\tfrac{1}{2}s-q+1)}\,f^{s+1}\!\int_{\kappa^{2}-f^{2}}^{\infty} t^{-q-\frac{1}{2}}(t+f^{2}-\kappa^{2})^{\frac{1}{2}s-q}(t+f^{2})^{-\frac{1}{2}s-q-1}\,dt,$}
\]
\[
\resizebox{\textwidth}{!}{$\displaystyle
\frac{\Gamma\tfrac{1}{2}s\,\Gamma\tfrac{1}{2}s}{\Gamma(\tfrac{1}{2}s+q)\,\Gamma(\tfrac{1}{2}-q)}\,f^{s+1}\!\int_{\kappa^{2}-f^{2}}^{\infty} t^{\frac{1}{2}s+q-1}(t+f^{2}-\kappa^{2})^{-q-\frac{1}{2}}(t+f^{2})^{-\frac{1}{2}s-q-\frac{1}{2}}\,dt,$}
\]
\[
\resizebox{\textwidth}{!}{$\displaystyle
\frac{\Gamma\tfrac{1}{2}s\,\Gamma\tfrac{1}{2}s}{\Gamma(\tfrac{1}{2}-q)\,\Gamma(\tfrac{1}{2}s+q)}\,f^{s+1}\!\int_{\kappa^{2}-f^{2}}^{\infty} t^{\frac{1}{2}s-q-1}(t+f^{2}-\kappa^{2})^{q-\frac{1}{2}}(t+f^{2})^{-\frac{1}{2}s+q-1}\,dt,$}
\]
\[
\resizebox{\textwidth}{!}{$\displaystyle
\frac{\Gamma\tfrac{1}{2}s\,\Gamma\tfrac{1}{2}s}{\Gamma(\tfrac{1}{2}s-q+1)\,\Gamma(\tfrac{1}{2}+q)}\,f^{s+1}\!\int_{\kappa^{2}-f^{2}}^{\infty} t^{-q-\frac{1}{2}}(t+f^{2}-\kappa^{2})^{\frac{1}{2}s-q-1}(t+f^{2})^{-\frac{1}{2}s+q-1}\,dt,$}
\]
where we may in the integrals write $t+\kappa^{2}-f^{2}$ in place of $t$, making the limits $\infty$, $0$; but the actual form is preferable.

\bigskip

\noindent\textbf{91.}\quad In the third form, for $f$ write $mf$, at the same time changing $t$ into $m^{2}t$; the new value of the disk-integral is
\[
\resizebox{\textwidth}{!}{$\displaystyle
= \frac{\Gamma\tfrac{1}{2}s\,\Gamma\tfrac{1}{2}s}{\Gamma(\tfrac{1}{2}-q)\,\Gamma(\tfrac{1}{2}s+q)}\,f^{s+1}\!\int_{m^{-2}(\kappa^{2}-f^{2})}^{\infty} t^{\frac{1}{2}s-q-1}\{m^{2}(t+f^{2})-\kappa^{2}\}^{q-\frac{1}{2}}(t+f^{2})^{-\frac{1}{2}s+q-1}\,dt.$}
\]
Writing here $mf=f+\delta f$, that is, $m=1+\dfrac{\delta f}{f}$, $m^{2}=1+\dfrac{2\delta f}{f}$, and observing that, if $-q+\tfrac{1}{2}$ be positive, the factor $\{m^{2}(t+f^{2})-\kappa^{2}\}^{-q+\frac{1}{2}}$ vanishes for the value $t=\dfrac{\kappa^{2}}{m^{2}}-f^{2}$ at the lower limit, we see that on this supposition, $-q+\tfrac{1}{2}$ positive, the value is
\[
\resizebox{\textwidth}{!}{$\displaystyle
= \frac{\Gamma\tfrac{1}{2}s\,\Gamma\tfrac{1}{2}s}{\Gamma(\tfrac{1}{2}-q)\,\Gamma(\tfrac{1}{2}s+q)}\,f^{s+1}\!\int_{\kappa^{2}-f^{2}}^{\infty} t^{\frac{1}{2}s-q-1}\!\left\{t+f^{2}-\kappa^{2} + \tfrac{2\delta f}{f}(t+f^{2})\right\}^{q-\frac{1}{2}}\!(t+f^{2})^{-\frac{1}{2}s+q-1}\,dt;$}
\]
viz.\ the term in $\delta f$ is $=\delta f$ multiplied by the expression
\[
\resizebox{\textwidth}{!}{$\displaystyle
2(-q+\tfrac{1}{2})\,\frac{\Gamma\tfrac{1}{2}s\,\Gamma\tfrac{1}{2}s}{\Gamma(\tfrac{1}{2}-q)\,\Gamma(\tfrac{1}{2}s+q)}\,f^{s}\!\int_{\kappa^{2}-f^{2}}^{\infty} t^{\frac{1}{2}s-q-1}(t+f^{2}-\kappa^{2})^{q-\frac{3}{2}}(t+f^{2})^{-\frac{1}{2}s+q}\,dt,$}
\]

\bigskip

\noindent\textbf{[p.~379]}\quad that is, multiplied by
\[
\resizebox{\textwidth}{!}{$\displaystyle
2\,(\tfrac{1}{2}-q)\,\frac{\Gamma\tfrac{1}{2}s\,\Gamma\tfrac{1}{2}s}{\Gamma(\tfrac{1}{2}+q)\,\Gamma(\tfrac{1}{2}s-q)}\,f^{s}\!\int_{\kappa^{2}-f^{2}}^{\infty} t^{\frac{1}{2}s-q-1}(t+f^{2}-\kappa^{2})^{q-\frac{3}{2}}(t+f^{2})^{-\frac{1}{2}s+q}\,dt,$}
\]
which is in fact $=\delta f$ multiplied by the value of the ring-integral.

\bigskip

\noindent\textbf{92.}\quad Comparing for the cases of an interior point $\kappa<f$ and an exterior point $\kappa>f$, the four expressions for the disk-integral, it will be noticed that only the third expressions correspond precisely to each other; viz.\ these are interior point, $\kappa<f$; the value is
\[
\resizebox{\textwidth}{!}{$\displaystyle
= \frac{\Gamma\tfrac{1}{2}s\,\Gamma\tfrac{1}{2}s}{\Gamma(\tfrac{1}{2}+q)\,\Gamma(\tfrac{1}{2}-q)}\,f^{s+1}\!\int_{0}^{\infty} t^{\frac{1}{2}s-q-1}(t+f^{2}-\kappa^{2})^{q-\frac{1}{2}}(t+f^{2})^{-\frac{1}{2}s+q-1}\,dt - \tfrac{\epsilon^{1-2q}}{\tfrac{1}{2}-q}\!\left(\frac{\Gamma\tfrac{1}{2}s\,\Gamma\tfrac{1}{2}}{\Gamma(\tfrac{1}{2}s+\tfrac{1}{2})}\right),$}
\]
where, if $\tfrac{1}{2}-q$ be positive (which is, in fact, a necessary condition in order to the applicability of the formula), the term in $\epsilon$ vanishes, and may therefore be omitted; and exterior point, $\kappa>f$; the value is
\[
\resizebox{\textwidth}{!}{$\displaystyle
= \frac{\Gamma\tfrac{1}{2}s\,\Gamma\tfrac{1}{2}s}{\Gamma(\tfrac{1}{2}+q)\,\Gamma(\tfrac{1}{2}-q)}\,f^{s+1}\!\int_{\kappa^{2}-f^{2}}^{\infty} t^{\frac{1}{2}s-q-1}(t+f^{2}-\kappa^{2})^{q-\frac{1}{2}}(t+f^{2})^{-\frac{1}{2}s+q-1}\,dt,$}
\]
differing only from the preceding one in the inferior limit $\kappa^{2}-f^{2}$ in place of $0$ of the integral. We have $\tfrac{1}{2}-q$ positive, and also $\tfrac{1}{2}s+q$ positive; viz.\ $q$ may have any value diminishing from $\tfrac{1}{2}$ to $-\tfrac{1}{2}s$, the extreme values not admissible.

\bigskip
\bigskip

\begin{center}
\textsc{Annex IV. Examples of Theorem A. Art.\ Nos.\ 93 to 112.}
\end{center}

\bigskip

\noindent\textbf{93.}\quad It is remarked in the text that, in the examples which relate to the $s$-coordinal sphere and ellipsoid respectively, we have a quantity $\theta$, a function of the coordinates $(a,\ldots,c,e)$ of the attracted point; viz.\ in the case of the sphere, writing $a^{2}+\ldots+c^{2}=\kappa^{2}$, we have
\[
\frac{\kappa^{2}}{f^{2}+\theta} + \frac{e^{2}}{\theta} = 1;
\]
in the case of the ellipsoid, we have
\[
\frac{a^{2}}{f^{2}+\theta} + \ldots + \frac{c^{2}}{h^{2}+\theta} + \frac{e^{2}}{\theta} = 1,
\]
the equations having in each case a positive root which is called $\theta$. The properties of the equation are the same in each case; but for the sphere, the equation being a quadric one, can be solved. The equation in fact is
\[
\theta^{2} - \theta(e^{2}+\kappa^{2}-f^{2}) - e^{2}f^{2} = 0,
\]
and the positive root is therefore
\[
\theta = \tfrac{1}{2}\!\left\{e^{2}+\kappa^{2}-f^{2} + \sqrt{(e^{2}+\kappa^{2}-f^{2})^{2}+4e^{2}f^{2}}\right\}.
\]

Suppose $e$ to diminish gradually and become $=0$; for an exterior point, $\kappa>f$, the value of the radical is $=\kappa^{2}-f^{2}$, and we have $\theta=\kappa^{2}-f^{2}$; for an interior point, $\kappa<f$,

\bigskip

\noindent\textbf{[p.~380]}\quad the value of the radical, supposing $e$ only indefinitely small, is $\displaystyle =f^{2}-\kappa^{2}+\frac{2e^{2}f^{2}}{f^{2}-\kappa^{2}}$, and we have $\displaystyle \theta=\tfrac{1}{2}\!\left(e^{2}+\kappa^{2}-f^{2}+f^{2}-\kappa^{2}+\frac{2e^{2}f^{2}}{f^{2}-\kappa^{2}}\right) = \frac{e^{2}f^{2}}{f^{2}-\kappa^{2}}$, or, what is the same thing, $\theta = e^{2}\!\left(1-\dfrac{\kappa^{2}}{f^{2}}\right)$; viz.\ the positive root of the equation continually diminishes with $e$, and becomes ultimately $=0$. If $\kappa$ or $e$ be indefinitely large, then the radical may be taken $=e^{2}+\kappa^{2}$, and we have $\theta$ indefinitely large, $=e^{2}+\kappa^{2}$.

\bigskip

\noindent\textbf{94.}\quad The result is similar for the general equation
\[
\frac{a^{2}}{f^{2}+\theta} + \ldots + \frac{c^{2}}{h^{2}+\theta} + \frac{e^{2}}{\theta} = 1;
\]
the left-hand side is $=0$ for $\theta=\infty$, and (as $\theta$ decreases) continually increases, becoming infinite for $\theta=0$; there is consequently a single positive value of $\theta$ for which the value is $=1$; viz.\ the equation has a single positive root, and $\theta$ is taken to denote this root.

In the last-mentioned equation, let $e$ gradually diminish and become $=0$; then for an exterior point, viz.\ if
\[
\frac{a^{2}}{f^{2}} + \ldots + \frac{c^{2}}{h^{2}} > 1,\quad\text{the equation } \frac{a^{2}}{f^{2}+\theta} + \ldots + \frac{c^{2}}{h^{2}+\theta} = 1
\]
has (as is at once seen) a single positive root, and $\theta$ becomes equal to the positive root of this equation; but for an interior point, or $\dfrac{a^{2}}{f^{2}} + \ldots + \dfrac{c^{2}}{h^{2}} < 1$, the equation just written down has no positive root, and $\theta$ becomes $=0$, that is, the positive root of the original equation continually diminishes with $e$, and for $e=0$ becomes ultimately $=0$; its value for $e$ small is, in fact, given by $\theta = e^{2}\!\left(1-\dfrac{a^{2}}{f^{2}}-\ldots-\dfrac{c^{2}}{h^{2}}\right)$. Also $a,\ldots,c,e$ (or any of them) indefinitely large, $\theta$ is indefinitely large, $=a^{2}+\ldots+c^{2}+e^{2}$.

\bigskip

\noindent\textbf{95.}\quad We have an interesting geometrical illustration in the case $s+1=2$; $\theta$ is here determined by the equation
\[
\frac{a^{2}}{f^{2}+\theta} + \frac{b^{2}}{g^{2}+\theta} + \frac{e^{2}}{\theta} = 1;
\]
viz.\ $\theta$ is the squared $a$-semiaxis of the ellipsoid, confocal with the conic $\dfrac{x^{2}}{f^{2}}+\dfrac{y^{2}}{g^{2}}=1$, which passes through the point $(a,b,e)$. Taking $e=0$, the point in question, if $\dfrac{a^{2}}{f^{2}}+\dfrac{b^{2}}{g^{2}}>1$, is a point in the plane of $xy$, outside the ellipse, and we have through the point a proper confocal ellipsoid, whose squared $\xi$-semiaxis does not vanish; but if $\dfrac{a^{2}}{f^{2}}+\dfrac{b^{2}}{g^{2}}<1$, then the point is within the ellipse, and the only confocal ellipsoid through the point is the indefinitely thin ellipsoid, squared semiaxes $(f^{2},g^{2},0)$, which in fact coincides with the ellipse.

\end{document}
